% GENERATED FILE. Edit canonical lessons, then run npm run generate:books.
% canonical-source-hash: fnv1a64:6dc2c50d903c7e2f
% canonical-lessons: LA-C444-malacia, LA-C444-tranquillitas, LA-C444-paeninsula, LA-C444-litoreus, LA-C444-enavigo

\chapter{Dead calm, Calm weather, Peninsula, Of the shore, To sail out}
\label{ch:la-dead-calm-calm-weather-peninsula-of-the-shore-to-sail-out}
\hlchaptermodality{444}

\begin{chapteropening}
\textbf{By the end of this chapter:} \emph{I can talk about calm seas and coasts: a dead calm, calm weather, a peninsula, things of the shore, and sailing out.}

Complete the last lesson of chapter 444: I can talk about calm seas and coasts: a dead calm, calm weather, a peninsula, things of the shore, and sailing out.
\end{chapteropening}

\section[malacia, malaciae]{malacia, malaciae — a dead calm (at sea)}

\label{lesson:LA-C444-malacia}



\begin{quote}
Before the new one: say the Latin for a strait, then the Latin for to sail past.
\end{quote}

\subsection*{You'll want to know: malacia, malaciae}
\textbf{malacia, malaciae} — \textquotedblleft{}a dead calm (at sea)\textquotedblright{}.

\begin{quote}
Malacia est; nāvēs in marī stant et nautae ventum exspectant.
\end{quote}

\begin{quote}
\emph{There is a dead calm; the ships stand still on the sea and the sailors wait for wind.}
\end{quote}

First declension, feminine: \emph{malacia, malaciae f.} No wind at all, the sailor's other trouble.

\begin{cousinweb}
\emph{malacia} was borrowed from Greek \emph{malakiā} 'softness, slackness' (EtymDB; Lewis \& Short), from \emph{malakos} 'soft': a sea gone slack.
\end{cousinweb}

\begin{grammarlens}[title={What you've built}]
One more word for telling, explaining and describing at B1.
\end{grammarlens}

\subsection*{Your turn}
\begin{itemize}
\raggedright
  \item \emph{Say it:} \emph{malacia, malaciae}
  \item \emph{Say it:} \emph{malacia, malaciae}, once more
  \item \emph{Say it:} a sentence of your own with \emph{malacia, malaciae}
  \item \emph{Recall:} say the Latin for a strait, then the Latin for to sail past, then say \emph{malacia, malaciae} again
\end{itemize}

\begin{tcolorbox}[breakable,colback=teal!4,colframe=teal!35!black,title={Before you move on}]
What does malacia, malaciae mean? (\textquotedblleft{}A dead calm (at sea)\textquotedblright{}.) Say it once more.
\end{tcolorbox}
\section[tranquillitās, tranquillitātis]{tranquillitās, tranquillitātis — calm weather, a calm; stillness}

\label{lesson:LA-C444-tranquillitas}



\begin{quote}
Before the new one: say the Latin for to sail past, then the Latin for a dead calm.
\end{quote}

\subsection*{You'll want to know: tranquillitās, tranquillitātis}
\textbf{tranquillitās, tranquillitātis} — \textquotedblleft{}calm weather, a calm; stillness\textquotedblright{}.

\begin{quote}
Post tempestātem tranquillitās venit.
\end{quote}

\begin{quote}
\emph{After the storm calm weather comes.}
\end{quote}

Third declension, feminine: \emph{tranquillitās, tranquillitātis f.} Also calm of mind. English \emph{tranquillity} comes from it.

\begin{cousinweb}
\emph{tranquillitās} is built on the adjective \emph{tranquillus} 'calm, still' with the suffix \emph{-itās} (EtymDB; Lewis \& Short). English \emph{tranquillity} comes from it.
\end{cousinweb}

\begin{grammarlens}[title={What you've built}]
One more word for telling, explaining and describing at B1.
\end{grammarlens}

\subsection*{Your turn}
\begin{itemize}
\raggedright
  \item \emph{Say it:} \emph{tranquillitās, tranquillitātis}
  \item \emph{Say it:} \emph{tranquillitās, tranquillitātis}, once more
  \item \emph{Say it:} a sentence of your own with \emph{tranquillitās, tranquillitātis}
  \item \emph{Recall:} say the Latin for to sail past, then the Latin for a dead calm, then say \emph{tranquillitās, tranquillitātis} again
\end{itemize}

\begin{tcolorbox}[breakable,colback=teal!4,colframe=teal!35!black,title={Before you move on}]
What does tranquillitās, tranquillitātis mean? (\textquotedblleft{}Calm weather, a calm; stillness\textquotedblright{}.) Say it once more.
\end{tcolorbox}
\section[paenīnsula, paenīnsulae]{paenīnsula, paenīnsulae — a peninsula}

\label{lesson:LA-C444-paeninsula}



\begin{quote}
Before the new one: say the Latin for a dead calm, then the Latin for calm weather.
\end{quote}

\subsection*{You'll want to know: paenīnsula, paenīnsulae}
\textbf{paenīnsula, paenīnsulae} — \textquotedblleft{}a peninsula\textquotedblright{}.

\begin{quote}
Oppidum in paenīnsulā est; mare ā tribus partibus est.
\end{quote}

\begin{quote}
\emph{The town is on a peninsula; the sea is on three sides.}
\end{quote}

First declension, feminine: \emph{paenīnsula, paenīnsulae f.}, also spelled \emph{pēnīnsula}. Almost an island. English \emph{peninsula} comes from it.

\begin{cousinweb}
\emph{paenīnsula} is built from \emph{paene} 'almost' and \emph{īnsula} 'island' (EtymDB; Lewis \& Short). English \emph{peninsula} comes from it.
\end{cousinweb}

\begin{grammarlens}[title={What you've built}]
One more word for telling, explaining and describing at B1.
\end{grammarlens}

\subsection*{Your turn}
\begin{itemize}
\raggedright
  \item \emph{Say it:} \emph{paenīnsula, paenīnsulae}
  \item \emph{Say it:} \emph{paenīnsula, paenīnsulae}, once more
  \item \emph{Say it:} a sentence of your own with \emph{paenīnsula, paenīnsulae}
  \item \emph{Recall:} say the Latin for a dead calm, then the Latin for calm weather, then say \emph{paenīnsula, paenīnsulae} again
\end{itemize}

\begin{tcolorbox}[breakable,colback=teal!4,colframe=teal!35!black,title={Before you move on}]
What does paenīnsula, paenīnsulae mean? (\textquotedblleft{}A peninsula\textquotedblright{}.) Say it once more.
\end{tcolorbox}
\section[lītoreus, lītorea, lītoreum]{lītoreus, lītorea, lītoreum — of the shore, on the beach}

\label{lesson:LA-C444-litoreus}



\begin{quote}
Before the new one: say the Latin for calm weather, then the Latin for a peninsula.
\end{quote}

\subsection*{You'll want to know: lītoreus, lītorea, lītoreum}
\textbf{lītoreus, lītorea, lītoreum} — \textquotedblleft{}of the shore, on the beach\textquotedblright{}.

\begin{quote}
Līberī harēnam lītoream in saccōs pōnunt.
\end{quote}

\begin{quote}
\emph{The children put beach sand into sacks.}
\end{quote}

An adjective of three endings, from \emph{lītus, lītoris} 'shore'. English \emph{littoral} comes from the related \emph{lītorālis}.

\begin{cousinweb}
\emph{lītoreus} is built on \emph{lītus, lītoris} 'shore' with the adjective suffix \emph{-eus} (EtymDB; Lewis \& Short). English \emph{littoral} comes from the related adjective \emph{lītorālis}.
\end{cousinweb}

\begin{grammarlens}[title={What you've built}]
One more word for telling, explaining and describing at B1.
\end{grammarlens}

\subsection*{Your turn}
\begin{itemize}
\raggedright
  \item \emph{Say it:} \emph{lītoreus, lītorea, lītoreum}
  \item \emph{Say it:} \emph{lītoreus, lītorea, lītoreum}, once more
  \item \emph{Say it:} a sentence of your own with \emph{lītoreus, lītorea, lītoreum}
  \item \emph{Recall:} say the Latin for calm weather, then the Latin for a peninsula, then say \emph{lītoreus, lītorea, lītoreum} again
\end{itemize}

\begin{tcolorbox}[breakable,colback=teal!4,colframe=teal!35!black,title={Before you move on}]
What does lītoreus, lītorea, lītoreum mean? (\textquotedblleft{}Of the shore, on the beach\textquotedblright{}.) Say it once more.
\end{tcolorbox}
\section[ēnāvigō, ēnāvigāre]{ēnāvigō, ēnāvigāre — to sail out, to sail away; to sail across}

\label{lesson:LA-C444-enavigo}



\begin{quote}
Before the new one: say the Latin for a peninsula, then the Latin for of the shore.
\end{quote}

\subsection*{You'll want to know: ēnāvigō, ēnāvigāre}
\textbf{ēnāvigō, ēnāvigāre} — \textquotedblleft{}to sail out, to sail away; to sail across\textquotedblright{}.

\begin{quote}
Nāvis ē portū ēnāvigat et in altum it.
\end{quote}

\begin{quote}
\emph{The ship sails out of the harbour and goes out to sea.}
\end{quote}

\emph{ēnāvigō, ēnāvigāre, ēnāvigāvī, ēnāvigātum}. With an accusative, to sail right across: \emph{mare ēnāvigāre}.

\begin{cousinweb}
\emph{ēnāvigō} is built from \emph{ē(x)-} 'out' and \emph{nāvigō} 'to sail' (Lewis \& Short).
\end{cousinweb}

\begin{grammarlens}[title={What you've built}]
One more word for telling, explaining and describing at B1.
\end{grammarlens}

\subsection*{Your turn}
\begin{itemize}
\raggedright
  \item \emph{Say it:} \emph{ēnāvigō, ēnāvigāre}
  \item \emph{Say it:} \emph{ēnāvigō, ēnāvigāre}, once more
  \item \emph{Say it:} a sentence of your own with \emph{ēnāvigō, ēnāvigāre}
  \item \emph{Recall:} say the Latin for a peninsula, then the Latin for of the shore, then say \emph{ēnāvigō, ēnāvigāre} again
\end{itemize}

\begin{tcolorbox}[breakable,colback=teal!4,colframe=teal!35!black,title={Before you move on}]
What does ēnāvigō, ēnāvigāre mean? (\textquotedblleft{}To sail out, to sail away; to sail across\textquotedblright{}.) Say it once more.
\end{tcolorbox}
